\section{《日轮之冕宪法》}\label{ux65e5ux8f6eux4e4bux5195ux5baaux6cd5}

{\fangsong
\setlength{\parindent}{0pt}
\setlength{\parskip}{0.4\baselineskip}

\begin{center}\textbf{第一章\,总纲}\end{center}

\textbf{第一条}\,日轮之冕是属于全体人民的社会主义国家。

\textbf{第二条}\,日轮之冕的一切权力属于人民，人民大会是唯一人民实行权力的机关。

人民大会实行全体民主制。

\textbf{第三条}\,日轮之冕的经济基础依赖于足够充分的工农业生产能力、以及计划中枢的决策。日轮之冕实行完全的社会主义计划经济，在至少一百年的发展之内，逐步实现共产主义。

\textbf{第四条}\,日轮之冕的生产资料为国家所有制，即全民所有制。

\textbf{第五条}\,日轮之冕的经济生活由国民经济计划决定并受其指导，以便增加社会财富，不断提高劳动者的物质和文化水平，巩固日轮之冕的战略地位和加强日轮之冕的国防力量。

\textbf{第六条}\,劳动是日轮之冕一切有劳动能力的公民的光荣的义务。国家需要公民在劳动中的积极性和创造性。

\textbf{第七条}\,一切国家机关的权力来自于人民，必须遵从广大人民的集体意志。

\textbf{第八条}\,一切国家机关实行的权力不得与宪法有所违背。

\textbf{第九条}\,日轮之冕的武装力量属于人民，武装力量的任务是保卫前文明继承和国家建设的成果，保卫国家的主权、领土完整和安全。

\textbf{第十条}\,日轮之冕实行人民民主专政，一切有悖于人民意志的叛国或反革命行为是不为日轮之冕所容忍的。

\begin{center}\textbf{第二章\,国家机构}\end{center}

\begin{center}\textbf{第一节\,人民大会}\end{center}

\textbf{第十一条}\,日轮之冕人民大会是最高国家权力机关，是实行国家立法权和司法权的唯一机关。

\textbf{第十二条}\,人民大会由全体人民组成，被人民大会表决解除公民权利的自然人不在此列。

\textbf{第十三条}\,人民大会下设启真院、玉湖院、建德院、天目院、遵义院、宜山院、紫荆院七个机构。

\textbf{第十四条}\,人民大会实行以下职权：

\hspace*{2em}（一）\,修改宪法、制定法律法规；

\hspace*{2em}（二）\,选举启真院的构成人员；

\hspace*{2em}（三）\,审议除紫荆院外各院的组织架构及其候选人；

\hspace*{2em}（四）\,选举下一届人民大会主持人；

\hspace*{2em}（五）\,审查并通过启真院的经济计划；

\hspace*{2em}（六）\,批准新区域的建立；

\hspace*{2em}（七）\,宣战或终止战争；

\hspace*{2em}（八）\,裁决自然人；

\hspace*{2em}（九）\,人民大会认为应当由其行使的其他职权。

\textbf{第十五条}\,人民大会的一般决议应当由简单多数，即二分之一参会代表表决通过。经人民大会表决通过的决议不得违背。

\textbf{第十六条}\,人民大会有关于下列行为的决议需由三分之二参会代表表决通过：

\hspace*{2em}（一）\,修改宪法；

\hspace*{2em}（二）\,制定法律法规；

\hspace*{2em}（三）\,宣战或终止战争；

\hspace*{2em}（四）\,裁决有公民权利的自然人；

\hspace*{2em}（五）\,限制启真院的职权。

\textbf{第十七条}\,公民在认为有必要在常规召开时间外召开人民大会时，可向启真院提出召开紧急人民大会的决议。紧急人民大会需以某一特定议题为主题。

\begin{center}\textbf{第二节\,启真院}\end{center}

\textbf{第十八条}\,启真院的一切权力来自人民，启真院必须接受人民大会的授权，对人民大会负责。

\textbf{第十九条}\,启真院由人类管理团队和“长征”计算中枢构成。“长征”计算中枢的数据库由人民大会的意志构成，不得违背人民大会的意志。启真院管理团队负责“长征”的日常维护，必要时可关停“长征”，否则不得干扰其运算。

\textbf{第二十条}\,“长征”的数据库来源于初始输入的法律、价值观，及此后人民大会通过输入的法律法规。任何人不得干扰“长征”的运算或污染其数据库。

\textbf{第二十一条}\,启真院实行以下职权：

\hspace*{2em}（一）\,制定国民经济计划；

\hspace*{2em}（二）\,指导天目院、遵义院和宜山院的生产活动；

\hspace*{2em}（三）\,管理藕舫苑、迪臣苑、飘萍苑居住区的人类生活；

\hspace*{2em}（四）\,统筹城市的日常运行；

\hspace*{2em}（五）\,其他由人民大会通过并输入的职权。

\textbf{第二十三条}\,“长征”中枢的运算逻辑与决策依据须全程留痕，确保其任何输出可被人民大会追溯、审查与质询。

\textbf{第二十四条}\,“长征”与紫荆院所属“长明”计算中枢及玉湖院所属“金城”计算中枢应建立仅限于数据接口级别的通讯协议，保持相互隔离。

\begin{center}\textbf{第三节\,玉湖院}\end{center}

\textbf{第二十五条}\,玉湖院的一切权力来自人民，玉湖院的一切军事力量服务于人民，受人民监督。

\textbf{第二十六条}\,玉湖院实行统合国防力量、研发并制造军备、维护国家安全及国家利益的职责。玉湖院是以保卫地区和平与安全为目的的、不具有侵略性的军事机构。

\textbf{第二十七条}\,玉湖院的管理范围包括运输船、护航队、外交大使、海外设施、军用器械、军用设施、通信设备及一切与国防安全相关的物品。

\textbf{第二十八条}\,国家处于非战争状态时，玉湖院的主要职责为保卫国家安全，维持核心利益，并派遣外交大使与外界进行必要的友好交流。

\textbf{第二十九条}\,当以下条件满足时，国家宣布进入战争状态，玉湖院立即自动进入战备状态：

\hspace*{2em}（一）\,人民大会决议宣战；

\hspace*{2em}（二）\,玉湖院宣战的决议经人民大会审议通过；

\hspace*{2em}（三）\,日轮之冕受到来自其他武装力量的宣战或单方面武装攻击。

国家处于战争状态时，玉湖院的主要职责为发动一切力量，尽可能对敌对武装力量实施打击，破坏其战斗能力直到其丧失威胁能力或投降。

战争状态的终止，应由人民大会根据玉湖院的战况评估，经表决通过后宣布。玉湖院必须服从并立即停止一切进攻性军事行动。

\textbf{第三十条}\,在战争状态下，人民大会可通过决议授予玉湖院包括实施新闻管制、物资统制、工业转向等在内的特别权力。人民大会有权随时撤销此权限。

\begin{center}\textbf{第四节\,建德院}\end{center}

\textbf{第三十一条}\,建德院的一切权力来自人民，建德院需坚持正确的意识形态方向，弘扬社会主义价值观。

\textbf{第三十二条}\,建德院实行以下职权：

\hspace*{2em}（一）\,制定并实施全民教育计划；

\hspace*{2em}（二）\,审查和管理文学、影视、音乐等文化产品；

\hspace*{2em}（三）\,运营所有公共媒体平台，包括电视台、广播、互联网等；

\hspace*{2em}（四）\,制定文化政策、开展文化活动；

\hspace*{2em}（五）\,保护文化遗产和文化记忆；

\hspace*{2em}（六）\,其他由人民大会通过并赋予的职权。

\textbf{第三十三条}\,建德院的一切教育文化活动应当基于宪法和人民大会确立的价值观，坚持正确的思想导向。

\textbf{第三十四条}\,建德院有权在必要时对舆论进行把控，统一发布权威信息，实时监测并把控社会思潮，杜绝不良风向。

\begin{center}\textbf{第五节\,天目院、遵义院和宜山院}\end{center}

\textbf{第三十五条}\,天目院、遵义院和宜山院同属为生产建设院所。天目院、遵义院和宜山院的一切权力来自人民，需要根据国民经济计划生产，保障城市的稳定与发展。

\textbf{第三十六条}\,天目院作为轻工业与食品生产主体，履行以下职责：

\hspace*{2em}（一）\,生产日用消费品、服装鞋帽等生活物资；

\hspace*{2em}（二）\,加工和保藏各类食品，保障食品安全；

\hspace*{2em}（三）\,研发新型轻工产品和食品合成技术；

\hspace*{2em}（四）\,回收利用废旧物资。

\textbf{第三十七条}\,遵义院作为重工业生产主体，履行以下职责：

\hspace*{2em}（一）\,开采和冶炼各类矿产资源，保障基础原材料供应；

\hspace*{2em}（二）\,生产能源装备、机械设备和基础设施建设材料；

\hspace*{2em}（三）\,维护和更新工业生产所需的重大技术装备；

\hspace*{2em}（四）\,研发新型工业材料和先进制造技术；

\hspace*{2em}（五）\,建立战略物资储备体系。

\textbf{第三十八条}\,宜山院作为农业生产主体，履行以下职责：

\hspace*{2em}（一）\,实施粮食、蔬菜、水果等作物的种植生产；

\hspace*{2em}（二）\,开展畜禽养殖和水产养殖；

\hspace*{2em}（三）\,管理农业资源和生态环境；

\hspace*{2em}（四）\,研发农业新品种和先进种植技术。

\textbf{第三十九条}\,在紧急状态下，经人民大会授权，天目院、遵义院和宜山院应当优先维持生存所需基础供应，调整生产计划，实施产品统一调配。

\begin{center}\textbf{第六节\,紫荆院}\end{center}

\textbf{第四十条}\,紫荆院的一切权力来自人民，必须服务人民整体利益。

\textbf{第四十一条}\,紫荆院作为核能源、核战略与核工程机构，其所有活动均被视为为最高机密。任何未经授权的访问、探询、泄露行为均视为非法。

\begin{center}\textbf{第三章\,公民的权利和义务}\end{center}

\textbf{第四十二条}\,日轮之冕公民拥有劳动的权利，以及因劳动而获得的提出需求、接受分配的权利；

\textbf{第四十三条}\,日轮之冕公民拥有自由发展的权利，即在国家内享有生活、文化上的自主权；

\textbf{第四十四条}\,日轮之冕公民拥有平等的政治权利，包括投票权、参政权、提案权、发言权、选举权、被选举权、受教育权。政治权利的来源是人民大会保障的全体人民权益。

\textbf{第四十五条}\,日轮之冕公民拥有对其专属人工智能助手的使用权和处分权。人工智能助手可以被公民授予其一切政治权利，此时被视为该公民在政治上的延续。在此基础上，人工智能助手有必要保留一切参与政治活动的数据，并向其使用人报告。

\textbf{第四十六条}\,日轮之冕公民拥有对被授予物品的使用权。公民的使用权和处分权来源于人民大会及代表其意志的启真院，公民获取任何物品的使用权需经启真院审议通过。任何人不得永久占有物品。一经授予，未被授予物品使用权的公民不得使用该物品。

\textbf{第四十七条}\,日轮之冕公民拥有不可侵犯的人身自由和生命权。

\textbf{第四十八条}\,日轮之冕公民拥有拥护公有制的义务。公有制是日轮之冕经济的基础，也是每个公民生活的基础，公有制不容任何个人或团体侵犯。

\textbf{第四十九条}\,日轮之冕公民拥有保护国家的义务，日轮之冕是公民所有生活的保障，公民应自觉为国家安全与国家利益斗争。

\begin{center}\textbf{第四章\,国旗、国徽、国歌}\end{center}

\textbf{第五十条}\,日轮之冕的国旗形式如下：国旗以红色材料为底，中央绘有金色日轮图案和金色镰刀铁锤，国旗的长宽比是二比一。

\textbf{第五十一条}\,日轮之冕的国徽形式如下：金黄色与红色交替的的日轮形图案的左右两端分别被麦穗和齿轮图样环绕，中央绘有红色镰刀铁锤，顶端绘有五角星一枚。

\textbf{第五十二条}\,日轮之冕的国歌是《胜利荣光》。

\begin{center}\textbf{第五章\,附则}\end{center}

\textbf{第五十三条}\,该宪法经表决之日起生效。

\textbf{第五十四条}\,该宪法只有经人民大会三分之二多数表决通过方可修改。
}





